% !TeX root = ../../tesis.tex

% ============================================================
%  §5.6.2 — El Espectro MJG: diagnóstico de la red actual
% ============================================================

\subsection{El Espectro MJG: diagnóstico de la red actual de la \zmvm}
\label{subsec:espectro-mjg}

% Valores confirmados (Notebook 06, 2026-09-18):
%   T   = 108.92 min  → T_norm = (180-108.92)/(180-85) = 0.748 → Aceptable
%   DI  = mediana ~1.4 → DI_norm = (2.5-1.4)/(2.5-1.0) = 0.733 → Aceptable
%   FC top = 92       → C_i: \pendiente{Q1/Q3 desde Notebook 05}
%   B(v) Tacubaya=0.2179 (top) → Gini: \pendiente{desde Notebook 06}
%   C   = \pendiente{C integrada ZMVM desde coverage_df, Notebook 04}

% =====================================================================
\subsubsection*{Apertura}
\label{subsubsec:espectro-apertura}
% =====================================================================

% §5.6.2.A — Párrafo de transición del instrumento a la aplicación real.
% Una oración que cierre el contexto de §5.6.1 y anuncie que los valores
% de las Fases 1–3 se clasifican a continuación.
% [POR REDACTAR — protocolo: 3 variantes → score → aprobación]

% =====================================================================
\subsubsection*{El Espectro MJG de la \zmvm{} 2026}
\label{subsubsec:espectro-tabla}
% =====================================================================

% §5.6.2.B — Tabla principal (tab:espectro-mjg) + párrafo de presentación.
% Columnas: Dimensión | Indicador | Valor bruto | Valor normalizado | Banda
% Párrafo anuncia qué muestra la tabla sin leer cada fila.
% [tab:espectro-mjg POR CREAR — requiere C y C_i_norm confirmados]

% =====================================================================
\subsubsection*{Lectura por dimensión}
\label{subsubsec:espectro-lectura}
% =====================================================================

% §5.6.2.C — Un solo párrafo integrado con las 5 dimensiones agrupadas
% por lo que revelan sobre el sistema:
%   • T + DI  → calidad de servicio una vez dentro de la red
%   • C       → alcance territorial
%   • C_i     → variedad de opciones nodales
%   • B(v)    → concentración de carga estructural
% [POR REDACTAR — requiere C y C_i_norm para D y B(v) Gini]

% =====================================================================
\subsubsection*{El patrón estructural de la \zmvm}
\label{subsubsec:espectro-sintesis}
% =====================================================================

% §5.6.2.D — Síntesis analítica: qué dice el perfil como conjunto.
% Hallazgo central del capítulo: la red llega al territorio pero
% concentra carga y muestra eficiencia en rango Aceptable, no Idóneo.
% [POR REDACTAR]

% =====================================================================
\subsubsection*{El Espectro MJG como línea base}
\label{subsubsec:espectro-cierre}
% =====================================================================

% §5.6.2.E — Cierre: este perfil establece la condición de referencia.
% Cap. 6 aplica el mismo instrumento a la red con el anillo y compara
% los dos perfiles dimensión a dimensión.
% [POR REDACTAR]
